\section{第\ref{sec:glutcomputer}章　\nt{GLUT}神经元计算什么}

本章的逻辑命题是关于非负权重电路的定理，在本章中推导得出；实验支持的是神经元模型，以及按这些定理搭建的电路（符合检测器、延迟线、自我维持的神经元群、链）确实存在于大脑之中。

{\footnotesize
\setlength{\tabcolsep}{3pt}\renewcommand{\arraystretch}{1.15}
\begin{longtable}{>{\raggedright}p{0.5cm}>{\raggedright}p{4.0cm}>{\raggedright}p{5.6cm}>{\raggedright}p{2.3cm}>{\raggedright\arraybackslash}p{1.7cm}}
\toprule
编号 & 命题 & 实验与测量内容 & 来源 & 状态 \\
\midrule\endfirsthead
\toprule
编号 & 命题 & 实验与测量内容 & 来源 & 状态 \\
\midrule\endhead
1 & 神经元是带阈值和复位的RC电路~\eqref{eq:glutneuron} & 向大鼠皮层锥体神经元胞体（脑片）注入类似活体突触流的噪声电流。平均频率随电流均值和散布变化的关系，可以用泄漏积分发放神经元加上慢的频率适应来描述。本章给出的 $\tau$ 和延迟范围没有引用实验 & \cite{Rauch2003} & 已测量 \\
2 & 模型~\eqref{eq:glutneuron} 是一种简化：输入分支本身会产生局部脉冲 & 在体直接记录小鼠视皮层锥体神经元的树突：视觉刺激引发依赖NMDA受体的局部树突脉冲；抑制它们会使神经元的朝向调谐变宽 & \cite{Smith2013} & 已测量 \\
3 & 符合窗口 $\Delta_{\max}=\tau\ln\frac{w}{\theta-w}$~\eqref{eq:window}；$\theta=1.5w$ 时等于 $0.7\tau$ & 模型~\eqref{eq:glutneuron} 的推论。对有快速抑制的神经元，有窄总和窗口的直接测量（第\ref{sec:gaba}章，其表中第~3行） & 本章推导 & 理论 \\
4 & 仅由\nt{GLUT}神经元组成的网络只能计算输入的单调函数（命题~\ref{prop:monotone}） & 本章的证明：从静默出发、右端不减的迭代。这是关于模型的命题；没有在无抑制的活体网络上检验它的实验 & 本章推导 & 理论 \\
5 & 感觉器官提供一对导线：一些细胞对刺激增强起反应，另一些对刺激减弱起反应 & 视网膜：视锥信号在双极细胞处就已分为接通和断开两个通道。在猴身上药理阻断接通型双极细胞：两个通道一直分开通到皮层；阻断后，检测光变强的能力变差，而检测光变弱的能力不受影响 & \cite{Schiller1992} & 已测量 \\
6 & 采用双线编码，\nt{GLUT}神经元网络可以异步地、无需共同节拍地计算任意逻辑函数，代价是神经元数目翻倍（命题~\ref{prop:dualrail}，\eqref{eq:dualrail}，图~\ref{fig:dualxor}） & 双线编码和凭信号本身判断就绪，是延迟不敏感异步电路的标准技术。六个神经元的异或电路在本章搭建。没有实验表明大脑正是用双线编码计算逻辑函数 & \cite{Sparso2001} & 理论 \\
7 & 人耳间声音到达时间差最大 $\approx0.6\ms$，而大脑能分辨约十微秒 & 二选一实验中的听者：时间差分辨阈（正确率75\,\%）对 $150$--$1700$\,Hz频带噪声为 $9\,\mu$s，对 $1$\,kHz纯音为 $11\,\mu$s，对咔嗒声为 $28\,\mu$s。上限 $0.6\ms$ 是本章的计算，$0.22\,\text{m}/343\,\text{m/s}$ & \cite{KlumppEady1956} & 已测量 \\
8 & 鸟类通过延迟线和符合检测器把时间差变成位置~\eqref{eq:itd}（图~\ref{fig:itd}） & 谷仓猫头鹰：来自两耳的轴突从相对两侧进入符合检测器核团；脉冲沿轴突的到达时间随深度有序变化，穿过核团的延迟范围（$700\um$ 上 $160\,\mu$s）与两耳间延迟范围一致 & \cite{Carr1990} & 已测量 \\
9 & 哺乳动物中没有找到延迟阵列；同样的任务由带有精确定时的\nt{GLYC}神经元抑制的符合检测器完成 & 沙鼠内侧上橄榄核在体记录：响应不形成方向图；用微量移液管施加甘氨酸及其阻断剂表明，精确定时的\emph{甘氨酸}抑制决定了延迟的工作范围。这里的抑制来自\nt{GLYC}，而不是\nt{GABA} & \cite{Brand2002,Grothe2010} & 已测量 \\
10 & 强反馈神经元群是触发器；自我维持的活动在刺激后持续数秒（图~\ref{fig:bistable}） & 猴前额叶皮层，延迟眼动任务，眼睛需移向记住的位置（延迟 $1$--$6$\,s）：$170$ 个任务相关神经元中有 $87$ 个在延迟期改变频率，其中 $79\,\%$ 依赖于记住的方向。这种活动由有两个稳定状态的反馈维持，是理论 & \cite{Funahashi1989,Wang2001} & 已测量 \\
11 & 驱动持续超过 $\approx0.7\tau$ 才能写入比特（图~\ref{fig:recurrentgroup}b） & 用欧拉法计算方程 $\tau\,\dd r/\dd t=-r+f(wr+I)$，$w=1$，$I=0.6$；\texttt{models/} 中没有对应脚本。没有实验 & 本章计算 & 模型 \\
12 & 记忆可以存放在突触易化中，几乎不需要脉冲 & 理论：末梢中的残余钙使记录保持约一秒。依据：雪貂内侧前额叶皮层（多个神经元同时记录）中存在一个由易化突触连接、后效持久的锥体神经元子网。关于记忆保持期间“安静”间歇的观察综述。皮层何时用哪种方式，是悬而未决的问题 & \cite{Mongillo2008,WangMarkram2006,Stokes2015} & 假说 \\
13 & 记忆由疲劳擦除：递质储备耗竭 & 皮层锥体神经元成对记录：高频脉冲下突触响应下降，下降速度由释放概率决定。这正是熄灭自我维持神经元群的机制，尚无直接证明 & \cite{Tsodyks1997} & 已测量 \\
14 & 神经元群链是由事件启动的定时器；鸣禽的神经元严格依次发放 & 在睡眠和鸣唱时记录斑胸草雀高级发声中枢中已鉴定的神经元：每个投射到运动核团的神经元在歌曲的一个精确时刻发出一个短簇放电，各神经元依次发放 & \cite{Hahnloser2002} & 已测量 \\
\bottomrule
\end{longtable}}
